% Generated by remark-latex
\label{sect:OIwikidocsindex}

\subsection*{欢迎来到 {\bfseries OI Wiki}！  }
\par \hyref{https://github.com/OI-wiki/OI-wiki}{\includegraphicsEverywhere{{OIwikidocsimageswordArt}}{Word Art}}\textsuperscript{\label{endnoteref:OIwikidocsindexOutLink_1-3}\hyperref[endnote:OIwikidocsindexOutLink_1]{[1]}}
\par {\bfseries OI}（Olympiad in Informatics，信息学奥林匹克竞赛）在中国起源于 1984 年，是五大高中学科竞赛之一．
\par {\bfseries ICPC}（International Collegiate Programming Contest，国际大学生程序设计竞赛）由 ICPC 基金会（ICPC Foundation）举办，是最具影响力的大学生计算机竞赛．由于以前 ACM 赞助这个竞赛，也有很多人习惯叫它 ACM 竞赛．
\par {\bfseries OI Wiki} 致力于成为一个免费开放且持续更新的 {\bfseries 编程竞赛（competitive programming）} 知识整合站点，大家可以在这里获取与竞赛相关的、有趣又实用的知识．我们为大家准备了竞赛中的基础知识、常见题型、解题思路以及常用工具等内容，帮助大家更快速深入地学习编程竞赛中涉及到的知识．
\par 本项目受 \hyref{https://ctf-wiki.org/}{​C​T​​F​​ ​​W​​i​​k​​i​}\textsuperscript{\label{endnoteref:OIwikidocsindexOutLink_2-1}\hyperref[endnote:OIwikidocsindexOutLink_2]{[2]}} 的启发，在编写过程中参考了诸多资料，在此一并致谢．


\subsubsection{参考资料与注释}
\begin{enumerate}
\renewcommand{\labelenumi}{[\theenumi]}
\item\label{endnote:OIwikidocsindexOutLink_1}\hyref{https://github.com/OI-wiki/OI-wiki}{} \hyperref[endnoteref:OIwikidocsindexOutLink_1-1]{[1-1]} \hyperref[endnoteref:OIwikidocsindexOutLink_1-2]{[1-2]} \hyperref[endnoteref:OIwikidocsindexOutLink_1-3]{[1-3]}\hfill \quad \qrcode[height=1cm]{https://github.com/OI-wiki/OI-wiki}
\renewcommand{\labelenumi}{\hyperref[endnoteref:OIwikidocsindexOutLink_2-1]{[\theenumi]}}
\item\label{endnote:OIwikidocsindexOutLink_2}\hyref{https://ctf-wiki.org/}{CTF Wiki}\hfill \quad \qrcode[height=1cm]{https://ctf-wiki.org/}
\end{enumerate}
